\documentclass[12pt,a4paper]{article}
\usepackage[T2A]{fontenc}
\usepackage[utf8]{inputenc}
\usepackage[russian]{babel}
\usepackage{amsmath,amssymb}
\usepackage[margin=2cm]{geometry}
\usepackage{listings}
\usepackage{xcolor}
\usepackage{hyperref}

\definecolor{titleblue}{RGB}{26,35,126}
\definecolor{codebg}{RGB}{30,30,30}
\definecolor{codefg}{RGB}{212,212,212}
\definecolor{kw}{RGB}{86,156,214}
\definecolor{fn}{RGB}{220,220,170}
\definecolor{cm}{RGB}{106,153,85}
\definecolor{st}{RGB}{206,145,120}

\lstdefinestyle{py}{
  language=Python,
  basicstyle=\ttfamily\small\color{codefg},
  backgroundcolor=\color{codebg},
  keywordstyle=\color{kw}\bfseries,
  commentstyle=\color{cm}\itshape,
  stringstyle=\color{st},
  functionnamestyle=\color{fn},
  numbers=none,
  frame=single,
  rulecolor=\color{codebg},
  breaklines=true,
  showstringspaces=false,
  tabsize=4,
  xleftmargin=1cm,
  xrightmargin=1cm,
  aboveskip=10pt,
  belowskip=10pt,
}

\lstset{style=py}

\title{\color{titleblue}\textbf{Описание функций сглаживания\\(Python, программистская версия)}}
\author{}
\date{}

\begin{document}
\maketitle

%% === 1. BOXCAR ===
\section{Boxcar}

\subsection{Сигнатура}

\begin{lstlisting}
def boxcar(win: int, data: list) -> list
\end{lstlisting}

\subsection{Параметры}

\begin{center}
\begin{tabular}{|l|l|l|}
\hline
\textbf{Параметр} & \textbf{Тип} & \textbf{Описание} \\
\hline
\texttt{win}  & \texttt{int}         & Размер окна усреднения \\
\texttt{data} & \texttt{list[float]} & Входной сигнал длиной $N$ \\
\hline
\end{tabular}
\end{center}

\subsection{Алгоритм}

\begin{lstlisting}
def boxcar(win, data):
    half = win // 2
    n = len(data)
    avg = []
    for i in range(n):
        start = max(0, i + 1 - half)
        end = min(half + i, n)
        sum_array = sum(data[start:end])
        aver = sum_array / len(data[start:end])
        avg.append(aver)
    return avg
\end{lstlisting}

Для каждого индекса \texttt{i} вычисляются границы окна \texttt{[start, end)} с обрезкой по краям массива. Вычисляется арифметическое среднее по срезу \texttt{data[start:end]}.

\subsection{Сложность}

\begin{center}
\begin{tabular}{|l|l|}
\hline
\textbf{Параметр} & \textbf{Значение} \\
\hline
Время            & $O(N \cdot \text{win})$ \\
Память           & $O(N)$ \\
Гранич. условия  & Clamping (обрезка окна) \\
\hline
\end{tabular}
\end{center}

%% === 2. GAUSSIAN ===
\section{Gaussian}

\subsection{Сигнатура}

\begin{lstlisting}
def gaussian(data: list, win: int) -> list
\end{lstlisting}

\subsection{Алгоритм --- ядро}

\begin{lstlisting}
def gaussian(data, win):
    half = win // 2
    sigma = win / 6.0
    n = len(data)
    k = [math.exp(-0.5*(i - half)**2 / sigma**2)
         for i in range(win)]
    k = [x/sum(k) for x in k]
\end{lstlisting}

\subsection{Алгоритм --- дополнение и свёртка}

\begin{lstlisting}
    p = ([data[i] for i in range(half-1, -1, -1)]
         + data
         + [data[i] for i in range(n-1, n-half-1, -1)])
    return [sum(k[j]*p[i+j] for j in range(win))
            for i in range(n)]
\end{lstlisting}

Ядро Гаусса: $w[k] = \exp\!\bigl(-0.5 \cdot (k - \text{half})^2 / \sigma^2\bigr)$, $\sigma = \text{win}/6$. Нормализация: $w\_norm[k] = w[k] / \sum w$. Дополнение: симметрическое отражение на \texttt{half} элементов. Свёртка: каждый выходной отсчёт --- взвешенная сумма по ядру.

\subsection{Сложность}

\begin{center}
\begin{tabular}{|l|l|}
\hline
\textbf{Параметр} & \textbf{Значение} \\
\hline
Время            & $O(N \cdot \text{win})$ \\
Память           & $O(N + \text{win})$ \\
Гранич. условия  & Mirror reflection (отражение) \\
\hline
\end{tabular}
\end{center}

%% === 3. PROGRESSIVE ===
\section{Progressive}

\subsection{Сигнатура}

\begin{lstlisting}
def progressive(data: list) -> list
\end{lstlisting}

\subsection{Алгоритм}

\begin{lstlisting}
def progressive(data):
    result = []
    total = 0.0
    for i, value in enumerate(data):
        total += value
        result.append(total / (i + 1))
    return result
\end{lstlisting}

Накапливается сумма \texttt{total}. На каждом шаге \texttt{i} вычисляется \texttt{total / (i + 1)} --- среднее отсчётов от $0$ до $i$.

\subsection{Сложность}

\begin{center}
\begin{tabular}{|l|l|}
\hline
\textbf{Параметр} & \textbf{Значение} \\
\hline
Время            & $O(N)$ \\
Память           & $O(N)$ \\
Гранич. условия  & Нет (окно растёт от $1$ до $N$) \\
\hline
\end{tabular}
\end{center}

%% === 4. HYBRID ===
\section{Hybrid}

\subsection{Сигнатура}

\begin{lstlisting}
def hybrid(data: list, win: int) -> list
\end{lstlisting}

\subsection{Алгоритм}

\begin{lstlisting}
def hybrid(data, win):
    n = len(data)
    box = boxcar(win, data)
    prog = progressive(data)

    out = []
    for i in range(n):
        if i < win:
            out.append(box[i])
        elif i < 2 * win:
            z = (i - win) / win
            k = (1.0 - z) * box[i] + z * prog[i]
            out.append(k)
        else:
            out.append(prog[i])
    return out
\end{lstlisting}

\subsection{Логика}

\begin{center}
\begin{tabular}{|l|l|}
\hline
\textbf{Условие} & \textbf{Действие} \\
\hline
\texttt{i < win}           & Берётся \texttt{boxcar[i]} \\
\texttt{win <= i < 2*win}  & Линейная интерполяция: $(1{-}z) \cdot box + z \cdot prog$ \\
\texttt{i >= 2*win}        & Берётся \texttt{progressive[i]} \\
\hline
\end{tabular}
\end{center}

Коэффициент интерполяции: $z = (i - \text{win}) / \text{win}$, $z \in [0, 1)$.

\subsection{Сложность}

\begin{center}
\begin{tabular}{|l|l|}
\hline
\textbf{Параметр} & \textbf{Значение} \\
\hline
Время            & $O(N \cdot \text{win})$ (из-за boxcar) \\
Память           & $O(N)$ \\
Гранич. условия  & Зависит от boxcar (clamping) \\
\hline
\end{tabular}
\end{center}

\end{document}
